\documentclass[11pt]{article}
\usepackage[T1]{fontenc}
\usepackage[utf8]{inputenc}
\usepackage{lmodern,amsmath,amssymb,amsthm,array,booktabs,longtable}
\usepackage[margin=1in]{geometry}
\usepackage[hidelinks]{hyperref}
\hypersetup{pdftitle={Explicit proofs of Mathar's recurrence and Barker's Hankel formulas from Mathar's and Cigler's theorems},pdfauthor={Haobo Ma; Wenlin Zhang},pdfsubject={Formal power series, recurrences and Hankel determinants}}
\emergencystretch=2em
\newtheorem{theorem}{Theorem}
\newtheorem{corollary}[theorem]{Corollary}
\newtheorem{lemma}[theorem]{Lemma}
\newcommand{\Q}{\mathbb{Q}}
\newcommand{\coeff}[2]{[x^{#1}]#2}
\title{Explicit proofs of Mathar's recurrence and Barker's Hankel formulas\\from Mathar's and Cigler's theorems}
\author{Haobo Ma\\\small ChronoAI Pte Ltd, Singapore
\and Wenlin Zhang\\\small National University of Singapore\\\small ChronoAI Pte Ltd, Singapore}
\date{16 November 2026}
\begin{document}
\maketitle
\begin{abstract}
Using Mathar's inverse-root theorem (2021, Theorem~2 and Example~4) and
Cigler's determinant theorem (2002 volume, Theorem~2), we give explicit
proofs of Mathar's three-term conjecture in OEIS A128057, revision~11, FORMULA
(2~December 2014), by coefficient pairing, and Barker's empirical Hankel
generating function in OEIS A128056, revision~5, FORMULA (27~June 2013),
by a block determinant, respectively.
For the coefficients of $(1+x)/\sqrt{1+4x^2}$ the three-term recurrence
holds for every $n\geq2$.
A polynomial identity gives the corresponding three-term recurrence for
$(u+vx)/\sqrt{1+cx^2}$ with arbitrary rational parameters.
The same determinant argument proves Barker's quotient generating function,
gives the unsigned and shifted Hankel transforms, and corrects related
sign, parity and indexing statements.
\end{abstract}

\section{The source question and the family}
Mathar's conjecture concerns the offset-zero sequence introduced by Paul
Barry in OEIS A128057~\cite{a128057}. Its initial terms are
$1,1,-2,-2,6,6,-20,-20,70,70$.
We prove the conjecture for every legal lag, $n\geq2$; the source itself
does not print an onset for the conjecture.
All power series below are formal series over $\Q$, and each square root
has constant coefficient one. Thus no analytic choice of branch is needed.

Mathar's general construction~\cite[Theorem~2 and Example~4]{mathar}
already treats this generating function and gives its four-term recurrence.
His Appendix~C explains how higher-order differential equations can reduce
the number of terms of a recurrence. The argument here makes the shorter
relation explicit by a parity calculation, and gives its uniform parameterized
form within Mathar's class. A direct recurrence identity below also makes
the conjecture an immediate consequence of his earlier theorem.
The classical central-binomial identities used in that calculation are
supporting identities, rather than separate novelty claims.
The paired coefficients also split the Hankel matrices into two moment
matrices. Their orthogonal Laurent-polynomial representation evaluates the
determinants and proves the empirical formulas of Barker~\cite{a128056,a128053}.
The determinant evaluation is a specialization of Cigler's Theorem~2
in the 2002 volume~\cite{cigler2002}, pp.~149--150, with weights
$(2,1,1,\ldots)$ and $z=1$; it is also covered by his later theorem
on polynomial moment combinations~\cite[Theorem~2, equations~(43)--(44)]{cigler2012}.
Callan's Theorem~1(i)~\cite{callan} supplies the general quadratic-root
coefficient formula, and Noe's equations~(2)--(4)~\cite{noe} supply the
coefficient and recurrence ingredients for the binomial-transform corollaries.
The earlier sources provide the general theorems; the recurrence combination
and parity-block specializations are displayed below.

\begin{theorem}\label{thm:family}
Let $u,v,c\in\Q$, let $S\in\Q[[x]]$ be the unique series satisfying
$S(0)=1$ and $S^2=1+cx^2$, and put
\[
 A(x)=(u+vx)S^{-1}=\sum_{n\geq0}a_nx^n,
 \qquad H=cu^2+v^2.
\]
Then, for $m\geq0$,
\begin{equation}\label{eq:pair}
 a_{2m}=u b_m,\qquad a_{2m+1}=v b_m,
 \qquad b_m=\left(-\frac c4\right)^m\binom{2m}{m}.
\end{equation}
For every integer $n\geq2$,
\begin{equation}\label{eq:general}
 n\bigl(Hn-2cu^2-v^2\bigr)a_n
 +cuv\,a_{n-1}
 +c(n-2)(Hn-cu^2)a_{n-2}=0.
\end{equation}
\end{theorem}

\begin{proof}
The factorial identity
$(m+1)\binom{2m+2}{m+1}=2(2m+1)\binom{2m}{m}$ gives
\begin{equation}\label{eq:adjacent}
 2(m+1)b_{m+1}+c(2m+1)b_m=0\qquad(m\geq0).
\end{equation}
Set $T(z)=\sum_{m\geq0}b_mz^m$. Equation~\eqref{eq:adjacent} is
$2(1+cz)T'(z)+cT(z)=0$. Hence
$((1+cz)T(z)^2)'=0$, and its constant coefficient is one, so
$(1+cz)T(z)^2=1$.
The unit $T(x^2)^{-1}$ is therefore a normalized square root of $1+cx^2$.
If $S_1$ and $S_2$ are two such roots, then
$(S_1-S_2)(S_1+S_2)=0$. The second factor has constant coefficient two
and is a unit, proving uniqueness. In particular $S^{-1}=T(x^2)$;
multiplication by $u+vx$ proves~\eqref{eq:pair}.

It remains to prove~\eqref{eq:general}. For $m\geq1$ define
\[
 K_m=2m b_m+c(2m-1)b_{m-1}=0,
\]
using~\eqref{eq:adjacent} at $m-1$.
Substitution of~\eqref{eq:pair} into the left side of~\eqref{eq:general}
at an even index $n=2m$ gives exactly
\[
 u\bigl((2m-2)cu^2+(2m-1)v^2\bigr)K_m.
\]
At an odd index $n=2m+1\geq3$ it gives exactly
\[
 v\bigl(2mcu^2+(2m+1)v^2\bigr)K_m.
\]
Both vanish. These two cases cover every $n\geq2$, including $n=2$.
\end{proof}

\begin{corollary}[Mathar's three-term recurrence]\label{cor:target}
The unique normalized square root of $1+4x^2$ has a two-sided inverse
with constant coefficient one. If
$(1+x)/\sqrt{1+4x^2}=\sum_{n\geq0}a_nx^n$, then
\begin{equation}\label{eq:target}
 n(5n-9)a_n+4a_{n-1}+4(5n-4)(n-2)a_{n-2}=0\quad(n\geq2),
\end{equation}
and $a_{2m}=a_{2m+1}=(-1)^m\binom{2m}{m}$.
Together with $a_0=a_1=1$, this recurrence determines the sequence uniquely.
\end{corollary}
\begin{proof}
Take $(u,v,c)=(1,1,4)$ in Theorem~\ref{thm:family}, so that $H=5$.
A unit in $\Q[[x]]$ has a unique two-sided inverse; its constant
coefficient here is one. Finally $n(5n-9)\ne0$ for every integer $n\geq2$,
which permits successive determination of $a_n$.
\end{proof}

\paragraph{A direct consequence of Mathar's theorem.}
Let $R_n$ be the left side of~\eqref{eq:target} and put
\[
 F_n=na_n+(n-2)a_{n-1}+4(n-1)a_{n-2}+4(n-3)a_{n-3}.
\]
Coefficient collection, for any sequence and $n\geq3$, gives
\begin{equation}\label{eq:bridge}
 R_n+R_{n-1}=(5n-9)F_n.
\end{equation}
Mathar's Theorem~2 and Example~4~\cite{mathar} give $F_n=0$ for
$n\geq4$; its printed bound is $n>\deg R=3$ in that theorem's notation.
The initial coefficients give $F_3=0$ and $R_2=0$.
Induction in~\eqref{eq:bridge} proves the conjecture for every $n\geq2$.
Thus the direct parity proof is an alternative derivation of a consequence
of the 2021 theorem, rather than an enlargement of its generating-function class.

\section{Other formulas in the source entries}
Differentiating $A=(u+vx)/\sqrt{1+cx^2}$ yields
\begin{equation}\label{eq:firstorder}
 (u+vx)(1+cx^2)A'+(cux-v)A=0.
\end{equation}
Indeed $\bigl(S^{-1}\bigr)'=-cxS^{-1}/(1+cx^2)$, so both terms cancel
after substitution. Extracting $x^{n-1}$ gives, for $n\geq3$,
\begin{equation}\label{eq:four}
 un a_n+v(n-2)a_{n-1}+cu(n-1)a_{n-2}+cv(n-3)a_{n-3}=0.
\end{equation}
The specialization $(1,1,4)$ is the four-term relation recorded by
Mathar in A128057 on 2~December 2014 and covered by his general
method~\cite{a128057,mathar}. This is distinct from~\eqref{eq:target}.

\paragraph{Unsigned sequence.}
Let $U_n=|a_n|$ for the sequence of Corollary~\ref{cor:target}.
Then $U_{2m}=U_{2m+1}=\binom{2m}{m}$, so
\[
 \sum_{n\geq0}U_nx^n=\frac{1+x}{\sqrt{1-4x^2}}.
\]
This proves the unsigned identification with A128014 and its
central-binomial description~\cite{a128014,a000984}. In particular it is
the parity formula printed in A128057: $(-1)^{\binom n2}=(-1)^m$
for both $n=2m$ and $n=2m+1$.
Theorem~\ref{thm:family} at $(1,1,-4)$ gives a three-term relation for
$U_n$; equation~\eqref{eq:four} at these parameters gives the four-term
relation recorded by Mathar in A128014 on 26~November 2012.

\begin{corollary}[Further unsigned identities]\label{cor:unsigned}
For $n\geq1$,
\[
 U_n=2^n\prod_{k=0}^{n-1}\left(\frac{k+1}{2n}\right)^{(-1)^k},
 \qquad
 \sum_{n\geq0}U_{n+1}\frac{x^n}{n!}=I_0(2x)+2I_1(2x),
\]
where $I_0(2x)=\sum_{m\geq0}x^{2m}/(m!)^2$ and
$I_1(2x)=\sum_{m\geq0}x^{2m+1}/(m!(m+1)!)$.
The integral formula in A128014 is valid as
\[
 U_n=\frac1\pi\operatorname{PV}\int_{-2}^2
 \frac{t^{n-1}(1+t)}{\sqrt{4-t^2}}\,dt \qquad(n\geq0).
\]
Only $n=0$ requires the Cauchy principal value at zero.
The number of binary abelian squares of length $n$ is $U_n$ for even $n$
and zero for odd $n$.
\end{corollary}
\begin{proof}
Separate the odd and even factors in the product. For $n=2m$ it gives
$\binom{2m}{m}$, and for $n=2m+1$ it gives the same value. Comparing
even and odd coefficients proves the exponential identity.
The substitution $t=2\cos\theta$ shows that the arcsine moments
are zero in odd degree and $\binom{2m}{m}$ in degree $2m$.
At degree $-1$ the symmetric principal value is zero. This proves the
integral statement, including $n=0$; the two one-sided integrals there
do not converge separately.
An abelian square consists of two equal-length words with the same number
of each symbol. For length $2m$ its count is
$\sum_{j=0}^m\binom mj^2=\binom{2m}{m}$ by the binomial theorem.
There is no such word of odd length. Thus A128014's unqualified
abelian-square identification needs the parity restriction: at $n=1$
the count is zero whereas $U_1=1$.
\end{proof}

\paragraph{Binomial transform.}
For $B_n=\sum_{k=0}^n\binom nk a_k$, the formal identity
$\sum_{n\geq k}\binom nk x^n=x^k/(1-x)^{k+1}$ gives
\begin{equation}\label{eq:binomial}
 B(x)=\sum_{n\geq0}B_nx^n
 =\frac1{1-x}A\!\left(\frac{x}{1-x}\right)
 =\frac1{(1-x)\sqrt{1-2x+5x^2}}.
\end{equation}
Thus this is the sequence A128058~\cite{a128058}.
Writing $D(x)=(1-2x+5x^2)^{-1/2}$ and expanding as
\[
 D(x)=\sum_{j\geq0}(-1)^j\binom{2j}{j}
             \frac{x^{2j}}{(1-x)^{2j+1}}
\]
gives $[x^k]D=\sum_{j=0}^{\lfloor k/2\rfloor}
(-1)^j\binom{k}{j}\binom{k-j}{j}$.
Since $B=D/(1-x)$, this also proves A128058's double-sum formula.
Differentiating~\eqref{eq:binomial} gives
\[
 (1-3x+7x^2-5x^3)B'-(2-8x+10x^2)B=0.
\]
Its coefficient of $x^{n-1}$ is
\begin{equation}\label{eq:binrec}
 nB_n+(1-3n)B_{n-1}+(7n-6)B_{n-2}+5(1-n)B_{n-3}=0
 \quad(n\geq3),
\end{equation}
reproving the relation recorded by Mathar on 23~January 2020.
It also follows directly from Noe's 2006 recurrence below by substituting
$d_n=B_n-B_{n-1}$.
Also $B_n$ is the partial sum of the coefficients of $D(x)$, as stated
in A128058. Those coefficients have the two sums and the recurrence
\[
 d_n=\sum_{j=0}^{\lfloor n/2\rfloor}(-1)^j\binom n{2j}\binom{2j}j,
 \quad
 nd_n=(2n-1)d_{n-1}-5(n-1)d_{n-2}\quad(n\geq2),
\]
with $d_0=d_1=1$. The sum equals the preceding expression for $[x^n]D$
term by term; the recurrence follows from
$(1-2x+5x^2)D'=(1-5x)D$. These reprove the corresponding formulas of
A098331, including Noe's recurrence (19~October 2005)~\cite{a098331}.

\section{Hankel determinants and Barker's formulas}
For a sequence $w$, define
$\mathcal H_n(w)=\det(w_{i+j})_{0\leq i,j\leq n}$.
Write $h_n=\mathcal H_n(a)$, $g_n=\mathcal H_n(U)$, and let
$L_0=2$, $L_1=1$, $L_{j+2}=L_{j+1}+L_j$ be the Lucas numbers.

\begin{lemma}\label{lem:moments}
For $b_j=(-1)^j\binom{2j}j$, put
$D_m=\det(b_{i+j})_{0\leq i,j<m}$ and
$E_m=\det(b_{i+j+1}-b_{i+j})_{0\leq i,j<m}$, with empty determinants one.
For $m\geq1$,
\[
 D_m=2^{m-1},\qquad E_m=(-1)^m2^{m-1}L_{2m}.
\]
For $b_j=\binom{2j}j$ instead, $D_m=2^{m-1}$ and
$E_m=2^m\cos(\pi m/3)$.
\end{lemma}
\begin{proof}
Let $\operatorname{CT}$ mean the coefficient of $z^0$ in a Laurent
polynomial and set $t=-2+z+z^{-1}$. Then
\[
 \operatorname{CT}t^j=(-1)^j\sum_{k=0}^j\binom jk^2
 =(-1)^j\binom{2j}j=b_j.
\]
Define monic polynomials $P_0=1$, $P_1(t)=t+2$,
$P_2(t)=(t+2)^2-2$, and
$P_m=(t+2)P_{m-1}-P_{m-2}$ for $m\geq3$.
After substitution, $P_m(t)=z^m+z^{-m}$ for $m\geq1$.
Consequently $\operatorname{CT}(P_iP_j)$ is zero for $i\ne j$,
one for $i=j=0$, and two for $i=j\geq1$.
Changing from monomials to these monic polynomials is triangular with
diagonal one, so $D_m=2^{m-1}$.

Orthogonality and monicity also give the determinant formula
\[
 D_m P_m(y)=\det\begin{pmatrix}
 b_0&b_1&\cdots&b_m\\
 b_1&b_2&\cdots&b_{m+1}\\
 \vdots&\vdots&&\vdots\\
 b_{m-1}&b_m&\cdots&b_{2m-1}\\
 1&y&\cdots&y^m
 \end{pmatrix}.
\]
Indeed the right side divided by $D_m$ is monic and orthogonal to
$1,y,\ldots,y^{m-1}$; the nonzero Gram determinant makes that polynomial
unique. At $y=1$, replace column $j$ by column $j$ minus column $j-1$,
from right to left. Expansion in the last row gives
$D_mP_m(1)=(-1)^mE_m$.
Now $P_1(1)=3$, $P_2(1)=7$, and
$P_m(1)=3P_{m-1}(1)-P_{m-2}(1)$ for $m\geq3$.
These are $L_{2m}$: the Lucas recurrence implies
$L_{2m}=3L_{2m-2}-L_{2m-4}$, with $L_2=3$, $L_4=7$.

For the unsigned case use $t=2+z+z^{-1}$ and replace $t+2$ by $t-2$
in the polynomials. The same proof gives $D_m=2^{m-1}$ and
$P_m(1)=2\cos(2\pi m/3)$, first for $m=1,2$ and then by the recurrence.
Thus $E_m=(-1)^mD_mP_m(1)=2^m\cos(\pi m/3)$.
\end{proof}

\begin{theorem}[The Hankel formulas]\label{thm:hankel}
We have $h_0=g_0=1$, and for $m\geq1$,
\begin{align}
 h_{2m-1}&=(-1)^m2^{2m-2}L_{2m},& h_{2m}&=2h_{2m-1},\label{eq:hankel}\\
 g_{2m-1}&=2^{2m-1}\cos(\pi m/3),& g_{2m}&=2g_{2m-1}.\label{eq:unsignedh}
\end{align}
The Hankel transform of $B$ equals $h$. In particular Barker's
empirical generating function in A128056 is valid:
\begin{equation}\label{eq:barker}
 \sum_{n\geq0}h_nx^n=
 -\frac{(2x-1)(4x^2-x+1)}{16x^4+12x^2+1}.
\end{equation}
\end{theorem}
\begin{proof}
Simultaneously reorder the rows and columns of the Hankel matrix by
even and odd indices; its determinant is unchanged. Since
$a_{2j}=a_{2j+1}=b_j$, replace the
row originally indexed $2i+1$ by that row minus row $2i$.
Do the analogous column subtraction. The off-diagonal blocks vanish.
For size $2m$ the diagonal blocks are
$(b_{i+j})_{i,j<m}$ and $(b_{i+j+1}-b_{i+j})_{i,j<m}$;
for size $2m+1$ they have sizes $m+1$ and $m$ respectively.
Therefore
$h_{2m-1}=D_mE_m$ and $h_{2m}=D_{m+1}E_m$.
Lemma~\ref{lem:moments} proves~\eqref{eq:hankel}; the same operations
with unsigned moments prove~\eqref{eq:unsignedh}.

The binomial invariance used here is Layman's theorem~\cite[Theorem~1]{layman},
whose matrix proof we include. Let $P=(\binom ij)_{0\leq i,j\leq n}$,
with $\binom ij=0$ for $j>i$.
Vandermonde's identity gives
$(B_{i+j})=P(a_{i+j})P^{\mathsf T}$: the coefficient of $a_k$ on the
right is $\sum_{r+s=k}\binom ir\binom js=\binom{i+j}k$.
Since $\det P=1$, the Hankel transforms agree.
Finally
\[
 \sum_{m\geq1}L_{2m}t^{m-1}=\frac{3-2t}{1-3t+t^2}.
\]
Using~\eqref{eq:hankel} at $t=-4x^2$ gives
$1-x(1+2x)(3+8x^2)/(1+12x^2+16x^4)$, which is~\eqref{eq:barker}.
\end{proof}

\begin{corollary}[The quotient and its correction]\label{cor:quotient}
Let $p_n=|g_n|$ and $q_n=h_n/p_n$. These are A128055 and A128053,
respectively. Put $\delta_m=2$ if $3\mid m$ and $\delta_m=1$ otherwise.
Then $q_0=p_0=1$, and for $m\geq1$,
\[
 p_{2m-1}=2^{2m-2}\delta_m,\quad p_{2m}=2^{2m-1}\delta_m,
 \qquad q_{2m-1}=q_{2m}=\frac{(-1)^mL_{2m}}{\delta_m}.
\]
In particular Barker's empirical formula in A128053 is valid:
\[
 \sum_{n\geq0}q_nx^n=
 -\frac{(x-1)(x^{10}-2x^9-5x^8+2x^7+9x^6+9x^4+2x^3-5x^2-2x+1)}
 {(x^4+3x^2+1)(x^8-3x^6+8x^4-3x^2+1)}.
\]
The trigonometric formula in that entry needs its factor
$(-1)^{\binom{n+1}2}$ to multiply both parity summands, not just the
even one. With this correction it is valid for every $n\geq0$,
using $F_{-1}=1$, $F_0=0$, $F_1=1$ for Fibonacci numbers.
\end{corollary}
\begin{proof}
For integer $m$, $2|\cos(\pi m/3)|=\delta_m$; hence all the formulas
follow from Theorem~\ref{thm:hankel}.
For the rational generating function, put $r_m=(-1)^mL_{2m}$.
Then $r_m=-3r_{m-1}-r_{m-2}$ and
$r_{3j}=-18r_{3j-3}-r_{3j-6}$ at admissible indices.
Thus the series of the paired values, including its constant term one, is
\[
 Q(t)=\frac{2+3t}{1+3t+t^2}-\frac{1+9t^3}{1+18t^3+t^6},
 \qquad
 \sum_{n\geq0}q_nx^n=1+\frac{1+x}{x}\bigl(Q(x^2)-1\bigr).
\]
Here $1/\delta_m=1-\tfrac12\boldsymbol1_{3\mid m}$, also for $m=0$;
this proves the first identity, and algebraic simplification proves the
displayed formula. For the trigonometric version,
$L_j=F_{j-1}+F_{j+1}$ and
$1/\delta_m=5/6-\cos(2\pi m/3)/3$.
The common sign is $(-1)^m$ for $n=2m-1,2m$.
At $n=1$ the formula as printed yields $+3$, whereas $q_1=-3$;
placing the sign outside both summands is the corrected formula.
The trigonometric formula in A128056 already has the common sign
$\cos(\pi n/2)-\sin(\pi n/2)$ and is correct.
\end{proof}

\begin{corollary}[The interleaved triangles]
In the named construction of A249095, interleave row $n$ of Pascal's
triangle with row $n-1$. Thus
\[
 T(n,2k)=\binom nk\ (0\leq k\leq n),\qquad
 T(n,2k+1)=\binom{n-1}k\ (0\leq k<n).
\]
Then $T(n,n)=U_n$, and scaling column $k$ by $2^k$, as in A249307,
gives diagonal $2^nU_n$, namely A249308~\cite{interleaved}.
\end{corollary}
\begin{proof}
Put $n=2m$ or $n=2m+1$ in the two formulas for $T$.
In each case the diagonal is $\binom{2m}m$.
The odd-column comment in A249095 prints $k-1$ in place of $k$;
at $(n,k)=(1,0)$ it would give zero instead of the middle entry one.
Each row has columns $0,\ldots,2n$, so the right edge of the scaled
triangle is at $2n$, not $2n+1$ as two A249307 lines print.
Its value is $2^{2n}=4^n$; the next column is outside the row.
\end{proof}

The exponents of $p_n$ are exactly A128054: writing $n=6j+s$,
$0\leq s<6$, they are $6j+e_s$, where
$(e_0,\ldots,e_5)=(0,0,1,2,3,5)$.
This lists multiples of six twice and omits $6j+4$, and proves
the exponent recurrence and generating function by summing six geometric
series and their derivatives. Likewise
\[
 \sum_{n\geq0}p_nx^n=\frac{1+x+2x^2+4x^3+8x^4+32x^5}{1-64x^6},
 \qquad \sum_{n\geq0}\frac1{p_n}=\frac{62}{21}.
\]
The latter follows by summing the first six reciprocals and multiplying
by $1/(1-1/64)$. This also proves the equivalent rational generating
function recorded in A128055~\cite{a128055,a128054}.
The alternating reciprocal sum of A128054 (Eldar, 2023) also follows.
If $k_n$ denotes that exponent sequence, cancellation of the two terms
with exponent $6j$ gives
\[
 \sum_{n\geq2}\frac{(-1)^n}{k_n}
 =\sum_{j\geq0}\left(\frac1{6j+1}-\frac1{6j+2}
                  +\frac1{6j+3}-\frac1{6j+5}\right)
 =\frac{5\pi}{12\sqrt3}-\frac{\log3}4+\frac{\log2}3.
\]
For the final equality, integrate the convergent geometric series on
$0\leq t<1$ to obtain the integrand
\[
 \frac{1-t+t^2-t^4}{1-t^6}
 =\frac1{3(1+t)}+\frac{1-t}{2(t^2+t+1)}
                   +\frac{t+1}{6(t^2-t+1)}.
\]
The grouped summands are $O((j+1)^{-2})$, so termwise integration is
valid. The two quadratic reciprocal integrals on $[0,1]$ are
$\pi/(3\sqrt3)$ and $2\pi/(3\sqrt3)$, respectively; their coefficients
are $3/4$ and $1/4$, and the logarithmic terms give the stated value.

For completeness, the shifted sequence has
\begin{equation}\label{eq:shifted}
 \mathcal H_n\bigl((U_{j+1})_{j\geq0}\bigr)
 =2^{n+1}\cos\!\left(\frac{\pi(n+1)}3\right),\qquad
 \sum_{n\geq0}\mathcal H_n\bigl((U_{j+1})_{j\geq0}\bigr)x^n
 =\frac{1-4x}{1-2x+4x^2}.
\end{equation}
Indeed $U_{j+1}=\operatorname{CT}(1+t)t^j$ with $t=z+z^{-1}$.
The same determinant formula at $y=-1$ gives, for size $N=n+1$,
$(-1)^N2^{N-1}\,2\cos(2\pi N/3)=2^N\cos(\pi N/3)$.
This proves the A128018 identification and Choulet's formula~\cite{a128018}.
For the original signed sequence with its first term removed, the same
method gives the earlier explicit result of Cigler~\cite[Theorem~2.2,
equation~(2.7)]{cigler2018}:
\begin{equation}\label{eq:signedshift}
 \mathcal H_n\bigl((a_{j+1})_{j\geq0}\bigr)
 =(-1)^{\binom{n+1}2}2^nL_{n+1}.
\end{equation}
To see this directly, write $t=i(z+z^{-1})$, where $i^2=-1$.
Then $a_{j+1}=\operatorname{CT}(1+t)t^j$.
The unweighted determinant of size $N$ is
$(-1)^{\binom N2}2^{N-1}$, and its monic orthogonal polynomial has
value $(-1)^NL_N$ at $-1$ for $N\geq1$: its first two values are
$-1,3$ and subsequent values satisfy $p_N=-p_{N-1}+p_{N-2}$.
The determinant formula proves~\eqref{eq:signedshift}.
For the unshifted sequence the Hankel values are
$1,1,2,-4,-8,\ldots$, as recorded in A128055, not
$1,2,2,1,1,\ldots$ as an A128017 comment states~\cite{a128017}.
The discrepancy already occurs at $n=1$, since
$\det\bigl(\begin{smallmatrix}1&1\\1&2\end{smallmatrix}\bigr)=1\ne2$.
Finally the Hankel transform of $d$, the A098331 coefficients above, is
$2^n(-1)^{\binom{n+1}2}$: its inverse-binomial-transform moments are
zero in odd degree and $(-1)^m\binom{2m}m$ in degree $2m$.
Scaling $t=z+z^{-1}$ by $i$ in the preceding orthogonal-polynomial
calculation multiplies the size-$N$ determinant by $i^{N(N-1)}$.
Binomial invariance then proves the formula.

\section{Coverage and verification}
\small
\begin{longtable}{@{}>{\raggedright\arraybackslash}p{0.35\textwidth}>{\raggedright\arraybackslash}p{0.61\textwidth}@{}}
\toprule
Source item & Coverage \\
\midrule
\endhead
A128057 three-term conjecture & Explicit consequence of Mathar, 2021: equation~\eqref{eq:bridge}; alternative parity proof, Corollary~\ref{cor:target}, every $n\geq2$. \\
Rational parameter family (not posed in the source) & Not posed; proved in Theorem~\ref{thm:family}, for every $u,v,c\in\Q$, including zero and degenerate parameters. \\
A128057 coefficient and transform comments & Reproved coefficient formula; prior general proof Callan, 2007, Theorem~1(i). Unsigned, binomial and Hankel consequences are explicit. \\
A128057 four-term relation & Reproved; prior relation Mathar, 2014; general proof Mathar, 2021, Theorem~2. \\
A128014 unsigned identification and four-term relation & Proved identification; reproved relation of Mathar, 2012. \\
A128014 product and shifted exponential formula & Proved: Corollary~\ref{cor:unsigned}; formulas recorded by Luschny, 2013, and Barry, 2010. \\
A128014 moment integral & Proved with Cauchy principal value at $n=0$; ordinary integral for $n\geq1$. \\
A128014 diagonal identifications with A249095 and A249308 & Proved for the named interleaving and column scaling. \\
A249095 odd-column index and A249307 row endpoint & Refuted-and-corrected: the interleaved-triangle corollary specifies the corrected indices and endpoint. \\
A128014 abelian-square comment & Refuted-and-corrected: parity restriction; counterexample $n=1$, Corollary~\ref{cor:unsigned}. \\
A128058 binomial transform and double sum & Proved: equation~\eqref{eq:binomial} and its expansion. \\
A128058 four-term relation & Consequence of Noe, 2006, equation~(4), by partial sums; later statement Mathar, 2020; equation~\eqref{eq:binrec}. \\
A128058 partial-sum and Hankel comments & Proved partial-sum identity; reproved binomial invariance, prior proof Layman, 2001, Theorem~1. \\
A128056 empirical generating function (Barker, 2013) & Explicit consequence of Cigler, 2002 volume, Theorem~2; later general proof Cigler, 2012, Theorem~2; block specialization in Theorem~\ref{thm:hankel}. \\
A128056 factorization and trigonometric formula & Proved: Theorem~\ref{thm:hankel}, Corollary~\ref{cor:quotient}; Fibonacci convention $F_{-1}=1$. \\
A128053 empirical generating function (Barker, 2013) & Proved: Corollary~\ref{cor:quotient}. \\
A128053 printed sign placement & Refuted-and-corrected: $q_1=-3$, not $+3$; common sign multiplies both summands. \\
A128055 signed-Hankel comment, generating function and reciprocal sum & Proved: equation~\eqref{eq:unsignedh} and six-residue formulas; prior statements Barry, 2008; Voznyy, 2009; Eldar, 2023. \\
A128054 listing, recurrence, generating function and reciprocal sum & Proved by six residue classes and the displayed integral; the printed trigonometric expression agrees in each class. \\
A128018 shifted-Hankel comment and cosine formula & Proved: equation~\eqref{eq:shifted}; prior statements Barry, 2009, and Choulet, 2008. \\
Shifted signed coefficients of A128057 & Reproved; prior explicit proof Cigler, 2018, Theorem~2.2(2.7); equation~\eqref{eq:signedshift}. \\
A128017 identification with the Hankel transform of A128014 & Refuted-and-corrected: $g_1=1$, not $2$; correct transform is~\eqref{eq:unsignedh}. \\
A098331 binomial, sum, recurrence and Hankel formulas & Reproved; prior recurrence Noe, 2005; Hankel statement Barry, 2009. \\
Other necklace/ballot interpretations in A128014 & Not covered. \\
Further identities of A098331, A128017 and A128018 & Not covered; no open-problem status is inferred. \\
\bottomrule
\end{longtable}
\normalsize

The declaration \texttt{result} in the
\href{https://github.com/the-omega-institute/trureturing/blob/2f8421f6b90d37b4815b887f6c60987553831b20/D5/S1/Recurrence/Algebraic/MatharA128057ThreeTermRecurrence.lean}{pinned \texttt{.lean} source}
states existence and uniqueness of the normalized square root of $1+4x^2$,
the normalized two-sided inverse, and equation~\eqref{eq:target} for every
natural $n\geq2$. The parameterized theorem and transform consequences
and Hankel results are proved in writing above. The supplied
\texttt{verify.py} gives exact symbolic and rational finite checks of
the recurrence, coefficient, transform and determinant identities;
its JSON output specifies the bounds. The finite checks complement
the universal proofs.

\begin{thebibliography}{99}
\bibitem{a128057}
P.~Barry and R.~J.~Mathar, contributions to OEIS A128057,
\emph{Expansion of $(1+x)/\sqrt{1+4x^2}$}, revision~11,
3~September 2021; recurrence conjecture dated 2~December 2014.
\url{https://oeis.org/A128057/internal}.
\bibitem{mathar}
R.~J.~Mathar, \emph{P-finite Recurrences From Generating Functions with
Roots of Polynomials}, arXiv:2109.02112v1 (2021),
Theorem~2, Example~4, and Appendix~C.
\url{https://arxiv.org/abs/2109.02112v1}.
\bibitem{layman}
J.~W.~Layman, \emph{The Hankel Transform and Some of its Properties},
Journal of Integer Sequences \textbf{4} (2001), Article 01.1.5,
Theorem~1, pp.~3--4; published 8~June 2001.
\url{https://cs.uwaterloo.ca/journals/JIS/VOL4/LAYMAN/hankel.pdf}.
\bibitem{cigler2002}
J.~Cigler, \emph{Some Relations Between Generalized Fibonacci and Catalan
Numbers}, Sitzungsberichte, Abteilung II \textbf{211} (2002 volume),
143--154, Theorem~2, equations~(2.1)--(2.5), proof on p.~150.
\url{https://homepage.univie.ac.at/johann.cigler/preprints/fib.cat.pdf}.
\bibitem{cigler2012}
J.~Cigler, \emph{Hankel determinants of some polynomial sequences},
arXiv:1211.0816v1, 5~November 2012, equations~(19)--(20),
Theorem~2, equations~(43)--(44).
\url{https://arxiv.org/abs/1211.0816v1}.
\bibitem{cigler2018}
J.~Cigler, \emph{Catalan numbers, Hankel determinants and Fibonacci polynomials},
arXiv:1801.05608v1, 17~January 2018, Theorem~2.2, equation~(2.7).
\url{https://arxiv.org/abs/1801.05608v1}.
\bibitem{callan}
D.~Callan, \emph{On Generating Functions Involving the Square Root of a
Quadratic Polynomial}, Journal of Integer Sequences \textbf{10} (2007),
Article 07.5.2, Theorem~1(i), pp.~2--4.
\url{https://cs.uwaterloo.ca/journals/JIS/VOL10/Callan/callan301.pdf}.
\bibitem{noe}
T.~D.~Noe, \emph{On the Divisibility of Generalized Central Trinomial
Coefficients}, Journal of Integer Sequences \textbf{9} (2006),
Article 06.2.7, equations~(2)--(4), Table~1.
\url{https://cs.uwaterloo.ca/journals/JIS/VOL9/Noe/noe35.pdf}.
\bibitem{a128014}
P.~Barry et al., OEIS A128014,
\emph{Central binomial coefficients $\binom{2k}{k}$ repeated},
FORMULA section, including Mathar's relation dated 26~November 2012.
\url{https://oeis.org/A128014}.
\bibitem{a128058}
P.~Barry and R.~J.~Mathar, contributions to OEIS A128058,
\emph{Expansion of $1/((1-x)\sqrt{1-2x+5x^2})$}, revision~7,
FORMULA section, recurrence dated 23~January 2020.
\url{https://oeis.org/A128058}.
\bibitem{a128056}
P.~Barry and C.~Barker, contributions to OEIS A128056,
\emph{Hankel transform of A128057}, revision~5,
empirical generating function dated 27~June 2013.
\url{https://oeis.org/A128056}.
\bibitem{a128053}
P.~Barry and C.~Barker, contributions to OEIS A128053,
\emph{Quotient of A128056 by A128055}, revision~5,
FORMULA; empirical generating function dated 27~June 2013.
\url{https://oeis.org/A128053}.
\bibitem{interleaved}
R.~Zumkeller et al., OEIS A249095 (revision~11), A249307 (revision~7)
and A249308 (revision~4), NAME and COMMENT sections.
\url{https://oeis.org/A249095}; \url{https://oeis.org/A249307};
\url{https://oeis.org/A249308}.
\bibitem{a128055}
P.~Barry et al., OEIS A128055, \emph{Powers of two with exponents A128054},
revision~16, COMMENT and FORMULA sections.
\url{https://oeis.org/A128055}.
\bibitem{a128054}
P.~Barry et al., OEIS A128054, \emph{Nonnegative integers with repeated
multiples of six and omitted integers congruent to four}, revision~21.
\url{https://oeis.org/A128054}.
\bibitem{a128017}
P.~Barry et al., OEIS A128017, revision~12, COMMENT section.
\url{https://oeis.org/A128017}.
\bibitem{a128018}
P.~Barry et al., OEIS A128018, revision~38, COMMENT and FORMULA sections;
cosine formula of R.~Choulet, 19~November 2008.
\url{https://oeis.org/A128018}.
\bibitem{a098331}
P.~Barry et al., OEIS A098331, revision~61, COMMENT and FORMULA sections;
recurrence of T.~D.~Noe, 19~October 2005.
\url{https://oeis.org/A098331}.
\bibitem{a000984}
OEIS Foundation Inc., OEIS A000984,
\emph{Central binomial coefficients: $\binom{2n}{n}=(2n)!/(n!)^2$}.
\url{https://oeis.org/A000984}.
\end{thebibliography}
\medskip
\noindent\small Copyright \textcopyright\ 2026 Haobo Ma and Wenlin Zhang.
This note is licensed under Creative Commons Attribution 4.0 International
(CC BY 4.0). Source entries are cited with their contributor attributions.
\end{document}
